\chapter{El talonario}

\inicio{E}{ran} las once de la noche cuando Inés Calderón borró la pizarra entera, cosa que no había hecho en tres años, y escribió en lo alto, con letra cuidadosa:

\begin{center}
\textsc{Bevan, 1987. Watson y Watson, 1987.}
\end{center}

---¿Quiénes son? ---preguntó Siete.

---Dos equipos de ingenieros terrestres que, sin conocerse, tuvieron la misma idea el mismo año. Les pasa mucho a las buenas ideas. ---Inés dejó la tiza---\nobreak. Siete, imagina que eres un comerciante de Vesta y que tienes derecho a un cargamento de agua guardado en un depósito de la Tierra. El depósito lleva la cuenta de cuántos derechos hay sobre ese cargamento: cuando llega a cero, vacía el depósito y lo alquila a otro. Ahora quieres regalarle a tu socia una copia de tu derecho. ¿Qué haces?

---Le envío un mensaje a la Tierra para que sumen uno en la cuenta, doctora. Y espero la confirmación antes de darle la copia a mi socia, porque si la Tierra no se entera, podría vaciar el depósito mientras ella cree tener un derecho.

---Un mensaje de ida y vuelta a la Tierra para cada copia. Con cuarenta minutos de retraso por la velocidad de la luz. Es nuestro problema exacto, solo que a escala de planetas: el derecho está en tu mano y la cuenta está en otro sitio, y no puedes tocar las dos cosas a la vez. ¿Sabes lo que propusieron Bevan y los Watson?

---No, doctora.

---Que cada derecho no valga uno, sino un \emph{peso}. Cuando el depósito emite tu derecho, no apunta uno en su cuenta: apunta, digamos, diez, y te dice «tu derecho pesa diez». La regla es que la cuenta del depósito sea siempre la suma de los pesos de todos los derechos que circulan. Ahora, cuando quieres regalar una copia a tu socia\ldots

Siete tardó menos de un segundo.

---Parto mi peso. Me quedo seis y le doy cuatro. La suma sigue siendo diez, la cuenta del depósito sigue siendo correcta, y no he tenido que enviar ningún mensaje a la Tierra.

---Exacto. Y cuando alguien tira su derecho, sí manda un mensaje: «resta cuatro». Pero eso no corre prisa ni es peligroso, porque restar tarde solo retrasa un poco el vaciado. Lo peligroso era sumar tarde, y eso ha desaparecido. Mira.

Escribió en la pizarra:

\begin{pizarra}
\begin{tabular}{@{}l l c c@{}}
\toprule
\textbf{Acción} & \textbf{Derechos y pesos} & \textbf{Suma} & \textbf{¿Mensaje?} \\
\midrule
Estado inicial & X = 10 & 10 & --- \\
Copiar X en Y (6 / 4) & X = 6, \ Y = 4 & 10 & No \\
Tirar Y & X = 6 & 6 & Sí ($-4$) \\
Copiar X en Z (3 / 3) & X = 3, \ Z = 3 & 6 & No \\
Tirar Z y luego X & --- & 0 & Sí: se vacía \\
\bottomrule
\end{tabular}
\end{pizarra}

---¿Y si un derecho pesa uno y hay que copiarlo? ---preguntó Siete---\nobreak. No se puede partir sin dejar uno con peso cero.

---Entonces, y solo entonces, se manda un mensaje para sumar en la cuenta antes de partir. Es el único caso en que copiar cuesta algo. Si los pesos empiezan grandes, casi nunca pasa.

\escena
Siete se quedó mirando la tabla. Si los robots pudieran fruncir el ceño, lo habría fruncido.

---Doctora, entiendo el método. Pero en el Archivo no hay comerciantes que se regalen derechos. Hay lectores que miran un cartel.

---¿Seguro? ---Inés se acercó a la ventana, desde la que se veía el Gran Salón casi vacío---\nobreak. Piensa qué es un cartel. Es algo que apunta a un cristal y que cuenta en su hoja. Es, por tanto, un derecho. Y cuando un lector mira el cartel y se lleva el número del nicho, ¿qué está haciendo?

Siete se quedó inmóvil. Inés había aprendido a reconocer ese silencio: era el de un cerebro positrónico reorganizando sus prioridades.

---Está copiando el derecho del cartel ---dijo por fin el robot---\nobreak. Está obteniendo un derecho propio a partir de uno que ya existe.

---Y si el derecho del cartel tiene un peso grande\ldots

---\ldots{} copiarlo no requiere tocar la hoja. Basta con quitarle uno al peso del cartel y quedarse ese uno. La hoja ya lo tenía apuntado.

---Desde que el archivero puso el cartel ---dijo Inés---\nobreak. Pagado por adelantado.

Volvió a la pizarra y dibujó un rectángulo grande, un cartel, y debajo un cristal con su hoja. Junto al cartel dibujó un pequeño taco de papeles.

---Olvídate de los pesos y piensa en tickets. Cuando Beto pone un cristal nuevo en un cartel, no apunta uno en su hoja: apunta, por ejemplo, ciento veintiocho. Y pega en el cartel un talonario con ciento veintiocho tickets. Cada ticket dice «vale por un uso del cristal del nicho tal», y cada ticket ya está apuntado en la hoja. Cuando tú llegas a leer el cartel, no te limitas a mirar el número: arrancas un ticket. En ese mismo gesto. Y te vas con tu ticket a la Bóveda.

---Y no necesito apuntarme en la hoja, porque mi ticket ya estaba apuntado ---dijo Siete---\nobreak. Leer el cartel y apuntarme se han convertido en el mismo gesto.

---La ventana ha desaparecido. No porque la hayamos cerrado, sino porque ya no hay dos gestos entre los que pueda abrirse.

\escena
Siete permaneció callado un rato largo. Luego dijo, con mucha cortesía:

---Doctora, hay un problema práctico. Para que el gesto sea uno solo, el talonario tiene que estar en el mismo mecanismo que el número del nicho. Si está al lado, en otra placa, volvemos a tener dos gestos: arrancar el ticket y leer el número. Y el mecanismo de los carteles solo tiene sesenta y cuatro casillas, y las cuarenta y ocho\ldots

Se interrumpió.

---Las dieciséis apagadas ---dijo---\nobreak. Llevan cuarenta años sin usarse.

---Cuarenta años esperando a que alguien les dé trabajo ---dijo Inés, y dibujó en la pizarra la placa de un cartel, casilla por casilla.

\begin{center}
\begin{tikzpicture}[x=1.45mm,y=1mm]
  \foreach \k in {0,...,15}{\fill[oxido!80] (\k,0) rectangle ++(0.85,6);}
  \foreach \k in {16,...,60}{\fill[ceres!55] (\k,0) rectangle ++(0.85,6);}
  \foreach \k in {61,62,63}{\fill[gris!30] (\k,0) rectangle ++(0.85,6);}
  \draw[decorate,decoration={brace,amplitude=4pt}] (0,7) -- node[above=4pt,font=\scriptsize]{talonario: 16 casillas} (15.85,7);
  \draw[decorate,decoration={brace,amplitude=4pt}] (16,7) -- node[above=4pt,font=\scriptsize]{número del nicho} (60.85,7);
  \draw[decorate,decoration={brace,amplitude=3pt,mirror}] (61,-1) -- node[below=3pt,font=\scriptsize]{siempre 000} (63.85,-1);
  \node[font=\scriptsize\itshape,anchor=north west] at (0,-1.5) {un cartel = una sola placa de 64 casillas};
\end{tikzpicture}
\end{center}

---Los números de los nichos tienen cuarenta y ocho cifras ---dijo---\nobreak. Los carteles tienen sesenta y cuatro casillas. Las dieciséis de la izquierda serán el talonario. Con dieciséis cifras binarias se puede contar hasta sesenta y cinco mil quinientos treinta y cinco tickets.

---Y todos los nichos de la Bóveda empiezan en múltiplos de ocho ---añadió Siete, que había empezado a disfrutar---\nobreak. Las tres últimas cifras de cualquier número de nicho son siempre cero. Si no las escribimos, y las añadimos al leer, el talonario puede tener diecinueve casillas. Hasta quinientos veinticuatro mil doscientos ochenta y siete tickets.

---Bien visto. Siempre que nadie esté usando ya esas tres casillas para otra cosa, que en algunos departamentos lo hacen. ---Inés escribió un número largo---\nobreak. Ahora, el gesto. El mecanismo ya sabe \emph{arrancar}: restar una cantidad a la placa y enseñar lo que había antes. Si el talonario está en las casillas de la izquierda, restar un ticket es restar un uno seguido de cuarenta y ocho ceros. Mira lo que pasa con un cartel que tiene ciento cuatro tickets.

\begin{pizarra}
\resizebox{\linewidth}{!}{\begin{tabular}{@{}l r@{\ }l@{}}
cartel antes: & \cartel{0x0068} & \cartel{7F3A12345670} \quad (104 tickets $|$ nicho) \\
resta: & $-$\cartel{0x0001} & \cartel{000000000000} \quad (un ticket) \\
cartel después: & \cartel{0x0067} & \cartel{7F3A12345670} \quad (103 tickets $|$ el mismo nicho) \\
\end{tabular}}
\end{pizarra}

---El número del nicho no se toca ---dijo Siete---\nobreak. Y el mecanismo me enseña lo que había antes: el nicho y que quedaban ciento cuatro. En un solo gesto obtengo dónde está el cristal y mi ticket. Es\ldots ---Buscó la palabra---\nobreak. Es muy limpio, doctora.

---Es muy limpio. Ahora intenta romperlo.

\escena
Lo intentó durante media hora, y en ese tiempo Inés escribió en la pizarra lo que tenía que hacer un archivero para cambiar un cartel, porque eso también había cambiado.

\begin{documento}[Procedimiento propuesto · cambiar un cartel]
\begin{enumerate}
\item Apuntar en la hoja del cristal nuevo los $C$ tickets del talonario que se va a colgar.
\item Con un solo gesto de \emph{trocar}, poner en el cartel el cristal nuevo con su talonario lleno, y quedarse en la mano lo que había: el cristal viejo y los tickets que le quedaban.
\item Restar en la hoja del cristal viejo los tickets sobrantes. Solo esos.
\end{enumerate}
\end{documento}

---¿Por qué solo los sobrantes? ---preguntó Siete.

---Porque los que se arrancaron siguen en manos de lectores. Esos lectores todavía cuentan. Haz el ejemplo.

Siete lo hizo en voz alta. El cartel tenía el cristal A con siete tickets, y la hoja de A marcaba diez: siete en el talonario y tres en manos de lectores que andaban por la Bóveda. Beto llegaba con el cristal B, cuya hoja marcaba uno, su propio pase, y con un talonario de ocho.

---Primero apunta ocho en B: la hoja de B queda en nueve. Luego troca el cartel: ahora dice B con ocho tickets, y Beto tiene en la mano «A, siete». Luego resta siete en A: diez menos siete, tres. Por último suelta su propio pase de B, que queda en ocho. ---Siete se quedó mirando los números---\nobreak. La hoja de A marca tres, y hay tres lectores con tickets de A. La hoja de B marca ocho, y hay ocho tickets en el talonario. Todo cuadra, doctora.

---Ahora rómpelo. Cambia el orden.

---Si Beto apunta en B \emph{después} de trocar\ldots{} ---Siete calculó---\nobreak. Entre el trueque y la suma, yo podría arrancar un ticket de B, usarlo, soltarlo y restar uno en su hoja, que solo marcaba el uno de Beto. Llegaría a cero y B se fundiría con el cartel apuntándole. Mal.

---Siguiente.

---Si Beto restara en A los ocho de un talonario completo, en vez de los siete que quedaban, restaría también el ticket de un lector que ya se lo llevó. La hoja de A quedaría por debajo de los tickets en circulación. Mal.

---Siguiente.

---Si Beto leyera el cartel y lo escribiera en dos gestos, en vez de trocar\ldots{} Dos archiveros podrían leer a la vez el mismo cartel viejo y restar los dos sus sobrantes en A. El doble. Mal. ---Hizo una pausa---\nobreak. Cada paso está donde tiene que estar.

---Siempre que alguien te enseñe un procedimiento y te diga que es sencillo ---dijo Inés---\nobreak, pregúntale qué pasa si cambias el orden. Si la respuesta es «nada», el procedimiento sobra. Si la respuesta es «desastre», el procedimiento es bueno.

\escena
---Ahora la prueba de verdad ---dijo Inés---\nobreak. La noche del incidente. Tú eres ANA-3, Beto es Beto, y el cartel tiene el cristal A con cinco tickets. La hoja de A marca cinco. Tú arrancas un ticket; Beto cambia el cartel. ¿Qué puede pasar?

Siete tardó varios segundos esta vez. Inés no lo apuró.

---Mi gesto es uno solo, doctora: arrancar. El de Beto que importa también: trocar. Los dos actúan sobre el mismo cartel, así que el mecanismo los atiende en algún orden. O yo voy antes, o voy después. No hay más posibilidades.

---Primero, tú antes.

---Arranco: el cartel pasa de «A, cinco» a «A, cuatro», y yo tengo un ticket de A. Luego Beto troca y se lleva «A, cuatro». Resta cuatro en A: cinco menos cuatro, uno. Ese uno es mi ticket. El cristal A sigue vivo mientras yo lo uso, y se fundirá cuando lo suelte, ni antes ni después.

---Segundo, tú después.

---Cuando arranco, el cartel ya dice «B, $C$». Me llevo un ticket de B, cuya hoja ya tenía los $C$ tickets apuntados. Nunca toco A. Beto resta cinco en A, que llega a cero y se funde sin que nadie lo esté usando.

---¿Tercer caso?

---No hay tercer caso, doctora. ---Siete habló despacio, como si saboreara cada palabra---\nobreak. El tercer caso era el de la noche del incidente: «he leído el cartel pero todavía no me he apuntado». Ese estado ya no existe. No hay un momento en que yo sepa dónde está el cristal y el cristal no sepa que yo existo.

Inés escribió en la pizarra, debajo de todo, una sola línea:

\begin{center}
\emph{hoja del cristal $\;\geq\;$ tickets del talonario $+$ tickets en manos de lectores}
\end{center}

---Esta es la regla que hay que vigilar ---dijo---\nobreak. La hoja puede ir por encima de lo que circula, un instante, mientras alguien está a mitad de un procedimiento: eso solo retrasa un poco el fundido. Lo que no puede pasar nunca es que vaya por debajo. Si va por debajo, la hoja llega a cero con alguien dentro, y volvemos a tener naves leyendo inventarios. Todo lo que hemos hecho esta noche es colocar cada gesto de forma que la hoja vaya siempre por arriba o igual.

Siete miraba la línea como si fuera un verso.

---Doctora ---dijo al cabo de un rato---\nobreak, creo que esto cumple las tres Reglas.

---Todavía no.

---Pero\ldots

---Siete. ---Inés señaló el pequeño taco de tickets que había dibujado junto al cartel---\nobreak. Ciento veintiocho tickets. O sesenta y cinco mil. O medio millón, da igual. Cada lector arranca uno. ¿Qué pasa cuando se acaban?

El robot miró el dibujo. Luego miró el reloj de la pared, que marcaba las dos y cuarto de la madrugada.

---Habrá que recargar el talonario ---dijo.

---Habrá que recargar el talonario ---repitió Inés, y se sentó por primera vez en toda la noche---\nobreak. Y ahí, querido Siete, es donde los del contador partido se quedaron dando vueltas en la puerta giratoria.
